\chapter*{Preface}
What is a number, a limit, a length? Every earlier volume used these words freely, and they worked. This book asks
what they really are, and the answers turn out to be among the most beautiful constructions in mathematics.

The real line comes first. It is built from nothing but sets: the natural numbers as the smallest set closed under
``add one'', the integers and the rationals as pairs, and the reals as limits of rationals, the one complete
ordered field there is. Cantor's diagonal argument then shows that there are more reals than rationals, and that
infinity comes in sizes.

Limits come next, freed from coordinates. A distance is enough to speak of convergence, completeness and fixed
points, and a single theorem, Banach's, solves equations, draws fractals and ranks web pages. Forgetting even the
distance leaves a topology, where continuity, connectedness and compactness make sense on their own. Compactness is
the art of treating infinite sets as if they were finite; Tychonoff's theorem, Arzel\`a--Ascoli and
Stone--Weierstrass are its masterpieces. Baire's theorem adds a surprise: in a complete space the ``typical'' object
is strange, and most continuous functions have a derivative nowhere. The same part ends with the Jordan curve theorem:
a closed curve that does not cross itself has an inside and an outside, proved from the winding number of a continuous curve and from the fixed-point theorem on a disc.

The second half of the book measures. Lebesgue's idea is to give a length to every reasonable set, and it cannot be
done for all of them: a set without a length exists, and a set can meet every interval in positive length while its complement does too.
From lengths comes the integral, built by slicing values rather than the axis, and with it limits pass under the
integral sign under conditions so mild that analysis and probability can be done with confidence. Volumes are
computed by slices, change by the determinant under a linear map, and give the volume of the ball in every
dimension, a number that tends to zero. Functions become points of spaces with a geometry, densities appear as
derivatives of one measure against another, and measures on metric spaces carry weak convergence and the dimension
of fractals, numbers like $\frac{\log8}{\log3}$, the dimension of the Sierpi\'nski carpet on the cover, strictly between one and two.

The goal is a stage on which modern analysis and probability can be performed with every step justified. Along the
way, results that were taken for granted earlier are proved at last: the intermediate value theorem, the existence of
the integral, the properties of $\exp$, $\cos$ and $\pi$, the existence of infinitely many independent coins, and
the theorem of Fubini. The beauty is in how little is needed: a few axioms about sets, one idea of distance, one idea
of measure, and an astonishing amount follows.

The book assumes the undergraduate analysis of Volumes~0 and~4, and some probability from Volume~3 for the examples; everything else is built here. The first five chapters are about sets, distances and topology, the last six about measure and integration, and each half can be read for its own sake. Volumes~7 to~11 stand on this ground: Hilbert spaces and Fourier analysis, harmonic analysis on groups, and probability as measure.

Each chapter ends with exercises, each with a hint and a complete solution.

The series asks whether pure order and pure chance can be told apart, and why the primes look random yet obey laws. This floor is the rooms: a limit, a length, Lebesgue's integral, and the Jordan curve, including its polygonal case. The next floor takes up the spectrum of a beam of infinite dimension. Conditional expectation, which the upper floors use as a beam, is a theorem here.
